%% Countably many dense requirements in a partially ordered set.

\subsection{The Rasiowa--Sikorski dense-set lemma}
\label{sec:p10-rasiowa-sikorski}

\begin{hbdefinition}{def:p10-dense-generic-filter}{Dense sets and generic filters in a poset}
Let \((P,\leq)\) be a nonempty partially ordered set.
The convention is that \(q\leq p\) means that \(q\) is a
\emph{stronger} condition than \(p\).
A subset \(D\subseteq P\) is \emph{downward dense} when
\[
 \forall p\in P\ \exists q\in D\quad q\leq p.
\]
A \emph{filter in the poset} is a nonempty \(G\subseteq P\) that
is upward closed and downward directed:
\[
\begin{aligned}
 q\in G,\ q\leq p&\quad\Longrightarrow\quad p\in G,\\
 p,q\in G&\quad\Longrightarrow\quad
       \exists r\in G\ (r\leq p\text{ and }r\leq q).
\end{aligned}
\]
For a family \(\mathcal D\) of downward dense subsets of \(P\),
the filter is \emph{\(\mathcal D\)-generic} when
\[
 G\cap D\ne\varnothing\qquad(D\in\mathcal D).
\]
Density requires a stronger member of \(D\) below each condition;
it imposes no closure condition on \(D\).
These are the order and genericity conventions of
\cite[\S8.1, pp.~130--132]{Ciesielski1997}.
\index{dense set!in a poset}\index{generic filter}\index{filter!poset}
\end{hbdefinition}

\begin{hbdefinition}{def:p10-dependent-choice}{Dependent choice for a prescribed start}
We use the following form of the axiom of dependent choice,
\(\mathsf{DC}\): if \(S\) is a nonempty set and
\(R\subseteq S\times S\) satisfies
\[
 \forall s\in S\ \exists t\in S\quad sRt,
\]
then for every prescribed \(s_0\in S\) there is a sequence
\((s_n)_{n\in\mathbb N}\) starting at that element and satisfying
\(s_nRs_{n+1}\) for every \(n\).
The next proof uses this principle to make successive choices of
stronger conditions, where each choice depends on the preceding one.
\index{dependent choice}
\end{hbdefinition}

\begin{hblemma}{thm:p10-rasiowa-sikorski}{proved=annex:p10.rs.1}{Rasiowa--Sikorski lemma}
Let \((P,\leq)\) be a nonempty poset, let \(p\in P\), and let
\(\mathcal D\) be a countable family of downward dense subsets of
\(P\), where countable includes finite.
Under \(\mathsf{ZF}+\mathsf{DC}\), there is a
\(\mathcal D\)-generic filter \(G\) containing \(p\).

If \(\mathcal D=\{D_n:n\in\mathbb N\}\) is given by a sequence,
one can take
\[
\begin{gathered}
 p_0=p,\qquad p_{n+1}\leq p_n,\qquad p_{n+1}\in D_n,\\
 G=\{q\in P:\exists n\in\mathbb N\quad p_n\leq q\}.
\end{gathered}
\]
If \(P\) has a fixed enumeration \(e:\mathbb N\to P\) onto it,
the same conclusion is provable in \(\mathsf{ZF}\) alone, by always
using the first enumerated suitable extension.
For a finite family \(\mathcal D\), \(\mathsf{ZF}\) suffices even
when \(P\) is arbitrary.  If \(P\) is finite, a filter containing
\(p\) can meet every downward dense subset of \(P\) at once in
\(\mathsf{ZF}\).
\end{hblemma}

The countable dense-set theorem and its descending-sequence construction
are \cite[Theorem~8.1.2, pp.~132--133]{Ciesielski1997}; the proof here
makes the sufficient choice assumptions explicit.
This theorem constructs a directed poset filter meeting specified dense
sets.  Theorem~\ref{thm:p10-prime-filter-separation} instead extends a
lattice filter to a prime filter while avoiding an ideal.
The arbitrary-poset conclusion above is proved using \(\mathsf{DC}\);
no implication from BPIT alone is asserted.
